\subsection{The actual algebraic representative in the threefold case}

\begin{theorem}[JAC-T4.4]
\label{thm:T44}
Let $A$ be an abelian threefold over an algebraically closed field of
arbitrary characteristic, let $h$ be an ample rational divisor class,
and accept 032, Corollary~8.4. Let $\operatorname{Ab}^2(A)$ be the
universal regular algebraic representative for algebraically trivial
codimension-two cycles. The algebraic projector onto
$H^3(A)_{\rm prim}$ induces an idempotent on
$\operatorname{Ab}^2(A)$ in the isogeny category. Write $B_{\rm reg}$
for its image. The canonical quotient homomorphism induces an isogeny
\[
 q:B_{\rm reg}\longrightarrow P^2(A)_{\rm prim}.
                                                        \tag{JAC.T44.1}
\]
The normalized primitive incidence divisor on $B_{\rm reg}$ is a
rational polarization; in the divisor convention of JAC-T4.3 it is
$\Delta^*d_{\rm reg}/2$. This assertion needs neither a universal
integral cycle nor surjectivity of the second Bloch map onto all of
$H^3(A)$.
\end{theorem}

\begin{proof}
We use the following precise known results. The codimension-two
algebraic representative exists, and a positive-degree miniversal
cycle exists for any surjective regular homomorphism:
Achter--Casalaina-Martin--Vial, \emph{Decomposition of the diagonal,
intermediate Jacobians, and universal codimension-2 cycles in positive
characteristic}, Ast\'erisque 455 (2025), \S1.1 and \S1.3, printed
pp.~3--6. Their Proposition~8.1, printed pp.~41--42, states that a
homologically zero correspondence induces zero on an algebraic
representative when the target codimension is $1$, $2$, or the
dimension. It is valid in every characteristic. We use only
codimension $2$.

Let $R=\operatorname{Ab}^2(A)$, with universal regular homomorphism
$\varphi:\operatorname{CH}^2(A)_{\rm alg}\to R(k)$.
Represent the primitive cohomological projector by a rational
correspondence $e$ supported in degree $3$. Proposition~8.1 applied
to $e^2-e$ shows that its action $e_R$ on $R$ is an actual rational
idempotent. Let $\iota:B_{\rm reg}\to R$ denote its inclusion.
Pull a miniversal family on $R\times A$ back by $\iota$ and divide
by its positive degree. Subtract its fibre over the origin. The
result is a rational family whose induced morphism is $\iota$.
Applying $e$ on the $A$ factor preserves this morphism.

We can furthermore arrange that the cohomology of this family is
entirely in $H^1(B_{\rm reg})\otimes H^3(A)_{\rm prim}$ without
altering its induced morphism. Here is an explicit justification
which does not assume that regular homomorphisms are incidence
homomorphisms. Put $b=\dim B_{\rm reg}$ and choose a rational
polynomial $Q$ with
\[
 Q(2)=1,\qquad Q(2^j)=0\quad(0\le j\le2b,\ j\ne1).
\]
Apply $Q(([2]\times1)^*)$ to the family. On cohomology this is the
degree-one projector on the parameter factor. On induced morphisms
to $R$ it acts as $Q(2)$, since the original morphism is a group
homomorphism and pullback by $[2]$ replaces it by $2\iota$.
Thus the induced morphism remains $\iota$. If $b=0$ the theorem is
immediate, so this construction may be made with $b>0$.
Denote the resulting normalized family by $z$.

Form
\[
 d_{\rm reg}=\operatorname{pr}_{12*}(z_{13}z_{23})_{11},
 \qquad \ell_{\rm reg}=\tfrac12\Delta^*d_{\rm reg}.
                                                        \tag{JAC.T44.2}
\]
For the normalization $f:C\to B_{\rm reg}$ of any integral curve,
translate so that a chosen point maps to zero and put
$z_C=(f\times1)^*z$. Its cohomological class is nonzero. Indeed,
if that class were zero, Proposition~8.1, applied with source $C$,
source codimension $1$ and target codimension $2$, would make the
induced homomorphism
$J(C)=\operatorname{Ab}^1(C)\to R$ zero. But this homomorphism
composed with the pointed Abel map is $\iota f$, which is
nonconstant. This is the pointed-family consequence of
Proposition~8.1; no assertion that an absolute constant family is
determined by its cohomology is needed.

JAC-T4.2 with $m=2$ now gives
\[
 \deg_C\ell_{\rm reg}
       =\tfrac12\deg_{C\times A}(z_C^2)>0.
                                                        \tag{JAC.T44.3}
\]
In particular $\ell_{\rm reg}$ is nef. It is also nondegenerate.
Otherwise the reduced identity component of the kernel of its
symmetric homomorphism would be a positive-dimensional abelian
subvariety $K$. The restriction of $\ell_{\rm reg}$ to $K$ has
zero associated homomorphism, hence is algebraically trivial. A
curve in $K$ would then have degree zero, contrary to
(JAC.T44.3). The nef, nondegenerate ampleness criterion used in
JAC-T4.3 therefore proves that $\ell_{\rm reg}$ is ample.

It remains to prove the comparison (JAC.T44.1), rather than assume
it. The incidence quotient map is a regular homomorphism. The
universal property of $R$ supplies a surjection
$R\to P^2(A)$. It commutes with correspondences, so its primitive
restriction $q$ is surjective in the isogeny category. Saito's
incidence identities, Proposition~4.4, identify
\[
 d_{\rm reg}=(q\times q)^*d_{P,\rm prim}
       \quad\hbox{in }\operatorname{NS}(B_{\rm reg}\times
                                      B_{\rm reg})_{\mathbf Q}.
                                                        \tag{JAC.T44.4}
\]
One can also check this identity directly: subtract the pullback
of a normalized Poincar\'e family on $P^2(A)$ from $z$. The
difference has zero incidence family map; testing it against either
family makes the mixed divisor component zero. Thus replacing one
family at a time leaves the $11$ incidence component unchanged.
If $q$ had positive-dimensional kernel, (JAC.T44.4) would make
$\ell_{\rm reg}$ numerically trivial on a curve in that kernel,
contradicting (JAC.T44.3). Its kernel is finite, proving the
isogeny and completing the proof.
\end{proof}

\paragraph{Comparison with existing autoduality results.}
The ACMV monograph proves substantially more integral information
under other hypotheses; those hypotheses must be kept separate.
Theorem~12.12(1), printed p.~68, assumes that the rational second
Bloch map is an isomorphism onto $H^3$. Its parts (2)--(3) use
universal $\operatorname{CH}_0$-triviality to construct a purely
inseparable isogeny or an isomorphism. Proposition~10.3, pp.~50--51,
establishes its standard torsion assumption in characteristic zero
or under the stated decomposition of the diagonal. Proposition~13.1
and Corollary~13.3, pp.~71--75, treat polarization by specialization.
None of these extra hypotheses was asserted here for an abelian
threefold. JAC-T4.4 instead uses 032 on products with curve Jacobians
and proves the rational \emph{primitive} assertion directly.

The monograph records its cup-product sign in Definition~12.1,
Remark~12.2 and Example~12.5, pp.~63--64. Its positivity assertion
also requires the appropriate primitive qualification or vanishing
of lower odd Hodge pieces. The raw full middle incidence form can
have both signs, as JAC-P03 shows. Comparisons of signed formulas
must use the specified divisor and alternating-pairing conventions;
JAC-T4.3 fixes ours by the elliptic-curve calculation.

\paragraph{Archive location.}
This strengthens the codimension-two interpretation of cote~44,
scan~2, Theorems~1--4, with the primitive qualification of scan~8.
It proves an isogeny comparison in this case; it does not identify
the incidence quotient with the universal regular representative in
every codimension or assert a principal polarization.
